\documentclass[12pt,aspectratio=169,ignorenonframetext]{ltxbeamer}

\usepackage[spanish,es-nodecimaldot,es-tabla]{babel}

\newcommand{\TikZ}{Ti\textit{k}Z}
\newcommand{\PGFPlots}{\textsc{pgfplots}}

\title{{\lmr\textmd{\PGFPlots}}}
\subtitle{Curso de {\lmr\LaTeX}}
\author{Joar Buitrago\texorpdfstring{\thanks{jebuitragoc@unal.edu.co}}{ <jebuitragoc@unal.edu.co>}}
\date{}

\sisetup{
    uncertainty-mode=separate,
}


\addbibresource{bibliography.bib}

\usepackage{pgfplots}
\usepackage{pgfplotstable}
\pgfplotsset{
    compat=1.18,
    legend style={
        fill=black,
        draw=white,
    },
}


\begin{document}
\begin{frame}
\titlepage
\end{frame}





\begin{frame}
\frametitle{Archivos fuente}

Vamos a trabajar con tres archivos

\begin{description}
    \item[battery.tsv] Archivo de datos en formato Tab Separated Values.
    \item[battery.csv] Archivo de datos en formato Comma Separated Values.
    \item[battery-title.csv] Archivo de datos en formato Comma Separated Values con títulos de tabla.
\end{description}
\end{frame}







\begin{frame}
\frametitle{Archivos fuente}

Los tres archivos contienen la misma información

\begin{center}
    \tiny
    \begin{tabular}{S[table-format=2.1(1)] S[table-format=0.6(1)]}
        \toprule
        {Resistencia [\unit{\ohm}]} & {Potencia [\unit{\watt}]} \\
        \midrule
        50.0 +- 2.5 & 0.036990 +- 0.007039 \\
        47.0 +- 2.5 & 0.036990 +- 0.007039 \\
        44.0 +- 2.5 & 0.038360 +- 0.007053 \\
        41.0 +- 2.5 & 0.042160 +- 0.007050 \\
        38.0 +- 2.5 & 0.043200 +- 0.007018 \\
        35.0 +- 2.5 & 0.046900 +- 0.007021 \\
        32.0 +- 2.5 & 0.048840 +- 0.006963 \\
        29.0 +- 2.5 & 0.052000 +- 0.006929 \\
        26.0 +- 2.5 & 0.055470 +- 0.006947 \\
        23.0 +- 2.5 & 0.060000 +- 0.006882 \\
        20.0 +- 2.5 & 0.064660 +- 0.006879 \\
        17.0 +- 2.5 & 0.069600 +- 0.006936 \\
        14.0 +- 2.5 & 0.079800 +- 0.007080 \\
        11.0 +- 2.5 & 0.086920 +- 0.007232 \\
         8.0 +- 2.5 & 0.085500 +- 0.007262 \\
         5.0 +- 2.5 & 0.101260 +- 0.008415 \\
         2.0 +- 2.5 & 0.091000 +- 0.011203 \\
         0.0 +- 2.5 & 0.081480 +- 0.011828 \\
        \bottomrule
    \end{tabular}
\end{center}
\end{frame}






\begin{frame}[fragile]
\begin{onlyenv}<1>
\begin{minted}{latex}
\begin{tikzpicture}
    \begin{axis}[
        width=0.75\paperwidth,
        height=0.75\paperheight,
        xlabel={Resistencia [\unit{\ohm}]},
        ylabel={Potencia [\unit{\watt}]},
    ]
        \addplot [
            red,
            mark=*,
            only marks,
        ] table {resources/battery.tsv};
    \end{axis}
\end{tikzpicture}
\end{minted}
\end{onlyenv}

\begin{onlyenv}<2->
\centering
\begin{tikzpicture}
    \begin{axis}[
        width=0.75\paperwidth,
        height=0.75\paperheight,
        xlabel={Resistencia [\unit{\ohm}]},
        ylabel={Potencia [\unit{\watt}]},
    ]
        \addplot [
            red,
            mark=*,
            only marks,
        ] table {resources/battery.tsv};
    \end{axis}
\end{tikzpicture}
\end{onlyenv}
\end{frame}





\begin{frame}[fragile]
\begin{onlyenv}<1>
\begin{minted}{latex}
\begin{tikzpicture}
    \begin{axis}[
        width=0.75\paperwidth,
        height=0.75\paperheight,
        xlabel={Resistencia [\unit{\ohm}]},
        ylabel={Potencia [\unit{\watt}]},
    ]
        \addplot [
            red,
            mark=*,
            only marks,
        ] table {resources/battery.csv};
    \end{axis}
\end{tikzpicture}
\end{minted}
\end{onlyenv}

\begin{onlyenv}<2->
\centering
\ttfamily
Package pgfplots Error: Sorry, the requested column number '1' in table 'resources/battery.csv' does not exist!? Please verify you used the correct index 0
\end{onlyenv}
\end{frame}





\begin{frame}[fragile]
\begin{minted}[escapeinside=||]{latex}
\usepackage{pgfplots}

\pgfplotstableset{|\textit{opciones}|}
\end{minted}
\end{frame}







\begin{frame}[fragile]
\begin{onlyenv}<1>
\begin{minted}{latex}
\pgfplotstableset{
    col sep=comma,
}
\end{minted}
\end{onlyenv}

\begin{onlyenv}<2>
\begin{minted}{latex}
\begin{tikzpicture}
    \begin{axis}[
        width=0.75\paperwidth,
        height=0.75\paperheight,
        xlabel={Resistencia [\unit{\ohm}]},
        ylabel={Potencia [\unit{\watt}]},
    ]
        \addplot [
            red,
            mark=*,
            only marks,
        ] table {resources/battery.csv};
    \end{axis}
\end{tikzpicture}
\end{minted}
\end{onlyenv}

\begin{onlyenv}<3->
\centering
\pgfplotstableset{
    col sep=comma,
}
\begin{tikzpicture}
    \begin{axis}[
        width=0.75\paperwidth,
        height=0.75\paperheight,
        xlabel={Resistencia [\unit{\ohm}]},
        ylabel={Potencia [\unit{\watt}]},
    ]
        \addplot [
            red,
            mark=*,
            only marks,
        ] table {resources/battery.csv};
    \end{axis}
\end{tikzpicture}
\end{onlyenv}
\end{frame}





\begin{frame}
\frametitle{Separadores de columna}

{\lmr\PGFPlots} posee configuración para diversos tipos de columnas.

\bigskip

\begingroup
\centering
\begin{tabular}{>{\ttfamily}l>{\ttfamily}lp{0.5\linewidth}}
    \toprule
    {\sffamily Opción} & {\sffamily RegExp} & Notas \\
    \midrule
    space & \textbackslash s+ & Al menos un espacio en blanco \\
    tab & \textbackslash t & \\
    comma & , & \\
    semicolon & ; & \\
    colon & : & \\
    braces & \{.*?\} & Todo el contenido dentro de las llaves se considera parte de la celda, incluyendo los saltos de línea \\
    \& & \& & \\
    ampersand & \& & \\
    \bottomrule
\end{tabular}\par
\endgroup
\end{frame}





\begin{frame}[fragile]
\begin{onlyenv}<1>
\begin{minted}{latex}
\begin{tikzpicture}
    \begin{axis}[
        width=0.75\paperwidth,
        height=0.75\paperheight,
        xlabel={Resistencia [\unit{\ohm}]},
        ylabel={Potencia [\unit{\watt}]},
    ]
        \addplot [
            red,
            mark=*,
            only marks,
        ] table [
            col sep=comma,
        ] {resources/battery.csv};
    \end{axis}
\end{tikzpicture}
\end{minted}
\end{onlyenv}

\begin{onlyenv}<2->
\centering
\begin{tikzpicture}
    \begin{axis}[
        width=0.75\paperwidth,
        height=0.75\paperheight,
        xlabel={Resistencia [\unit{\ohm}]},
        ylabel={Potencia [\unit{\watt}]},
    ]
        \addplot [
            red,
            mark=*,
            only marks,
        ] table [
            col sep=comma,
        ] {resources/battery.csv};
    \end{axis}
\end{tikzpicture}
\end{onlyenv}
\end{frame}






\begin{frame}[fragile]
\begin{onlyenv}<1>
\begin{minted}{latex}
\begin{tikzpicture}
    \begin{axis}[...]
        \addplot [...] table [
            x index=1,
            y index=0,
        ] {resources/battery.csv};
    \end{axis}
\end{tikzpicture}
\end{minted}
\end{onlyenv}

\begin{onlyenv}<2->
\centering
\pgfplotstableset{
    col sep=comma,
}
\begin{tikzpicture}
    \begin{axis}[
        width=0.75\paperwidth,
        height=0.75\paperheight,
        ylabel={Resistencia [\unit{\ohm}]},
        xlabel={Potencia [\unit{\watt}]},
    ]
        \addplot [
            red,
            mark=*,
            only marks,
        ] table [
            x index=1,
            y index=0,
        ] {resources/battery.csv};
    \end{axis}
\end{tikzpicture}
\end{onlyenv}
\end{frame}





\begin{frame}[fragile]
\begin{onlyenv}<1>
\begin{minted}{latex}
\begin{tikzpicture}
    \begin{axis}[...]
        \addplot [...] table [
            x index=1,
            y=Resistance,
        ] {resources/battery-title.csv};
    \end{axis}
\end{tikzpicture}
\end{minted}
\end{onlyenv}

\begin{onlyenv}<2->
\centering
\pgfplotstableset{
    col sep=comma,
}
\begin{tikzpicture}
    \begin{axis}[
        width=0.75\paperwidth,
        height=0.75\paperheight,
        ylabel={Resistencia [\unit{\ohm}]},
        xlabel={Potencia [\unit{\watt}]},
    ]
        \addplot [
            red,
            mark=*,
            only marks,
        ] table [
            x index=1,
            y=Resistance,
        ] {resources/battery-title.csv};
    \end{axis}
\end{tikzpicture}
\end{onlyenv}
\end{frame}






\begin{frame}[fragile]
\begin{onlyenv}<1>
\begin{minted}{latex}
\begin{tikzpicture}
    \begin{axis}[...]
        \addplot [...] table [
            x index=1,
            y expr=123,
        ] {resources/battery-title.csv};
    \end{axis}
\end{tikzpicture}
\end{minted}
\end{onlyenv}

\begin{onlyenv}<2->
\centering
\pgfplotstableset{
    col sep=comma,
}
\begin{tikzpicture}
    \begin{axis}[
        width=0.75\paperwidth,
        height=0.75\paperheight,
        ylabel={Resistencia [\unit{\ohm}]},
        xlabel={Potencia [\unit{\watt}]},
    ]
        \addplot [
            red,
            mark=*,
            only marks,
        ] table [
            x index=1,
            y expr=123,
        ] {resources/battery-title.csv};
    \end{axis}
\end{tikzpicture}
\end{onlyenv}
\end{frame}






\begin{frame}[fragile]
\frametitle{Expresiones}

Las expresiones para columnas de {\lmr\PGFPlots} pueden incluir los siguientes comandos.

\bigskip

\begingroup
\centering
\begin{tabular}{lp{0.5\linewidth}}
    \toprule
    Comando & Detalle \\
    \midrule
    \mintinline[escapeinside=||]{latex}{\thisrow{|\textit{columna}|}} & Valor de la fila actual en la columna con el nombre \texttt{\textit{columna}} \\
    \mintinline[escapeinside=||]{latex}{\thisrowno{|\textit{columna}|}} & Valor de la fila actual en la columna con el índice \texttt{\textit{columna}} \\
    \mintinline{latex}{\coordindex} & Índice de la fila actual \\
    \mintinline{latex}{\lineno} & Línea actual del documento \\
    \bottomrule
\end{tabular}\par
\endgroup
\end{frame}






\begin{frame}[fragile]
\begin{onlyenv}<1>
\begin{minted}{latex}
\begin{tikzpicture}
    \begin{axis}[...]
        \addplot [...] table [
            y expr=100 * \thisrowno{1},
        ] {resources/battery-title.csv};
    \end{axis}
\end{tikzpicture}
\end{minted}
\end{onlyenv}

\begin{onlyenv}<2->
\centering
\pgfplotstableset{
    col sep=comma,
}
\begin{tikzpicture}
    \begin{axis}[
        width=0.75\paperwidth,
        height=0.75\paperheight,
        xlabel={Resistencia [\unit{\ohm}]},
        ylabel={Potencia [\unit{\watt}]},
    ]
        \addplot [
            red,
            mark=*,
            only marks,
        ] table [
            y expr=100 * \thisrowno{1},
        ] {resources/battery-title.csv};
    \end{axis}
\end{tikzpicture}
\end{onlyenv}
\end{frame}






\begin{frame}[fragile]
\begin{onlyenv}<1>
\begin{minted}{latex}
\begin{tikzpicture}
    \begin{axis}[...]
        \addplot [...] table [
            y expr=sqrt(\thisrow{Power}),
        ] {resources/battery-title.csv};
    \end{axis}
\end{tikzpicture}
\end{minted}
\end{onlyenv}

\begin{onlyenv}<2->
\centering
\pgfplotstableset{
    col sep=comma,
}
\begin{tikzpicture}
    \begin{axis}[
        width=0.75\paperwidth,
        height=0.75\paperheight,
        xlabel={Resistencia [\unit{\ohm}]},
        ylabel={Potencia [\unit{\watt}]},
    ]
        \addplot [
            red,
            mark=*,
            only marks,
        ] table [
            y expr=sqrt(\thisrow{Power}),
        ] {resources/battery-title.csv};
    \end{axis}
\end{tikzpicture}
\end{onlyenv}
\end{frame}






\begin{frame}[fragile]
\begin{onlyenv}<1>
\begin{minted}{latex}
\begin{tikzpicture}
    \begin{axis}[...]
        \addplot [
            ...,
            error bars,
            x dir=both,
            x explicit,
            y dir=both,
            y explicit,
        ] table [
            x error=ResistanceError,
            y error=PowerError,
        ] {resources/battery-title.csv};
    \end{axis}
\end{tikzpicture}
\end{minted}
\end{onlyenv}

\begin{onlyenv}<2->
\centering
\pgfplotstableset{
    col sep=comma,
}
\begin{tikzpicture}
    \begin{axis}[
        width=0.75\paperwidth,
        height=0.75\paperheight,
        xlabel={Resistencia [\unit{\ohm}]},
        ylabel={Potencia [\unit{\watt}]},
    ]
        \addplot [
            red,
            mark=*,
            only marks,
            error bars,
            x dir=both,
            x explicit,
            y dir=both,
            y explicit,
        ] table [
            x error=ResistanceError,
            y error=PowerError,
        ] {resources/battery-title.csv};
    \end{axis}
\end{tikzpicture}
\end{onlyenv}
\end{frame}





\begin{frame}[fragile]
\begin{onlyenv}<1>
\begin{minted}{latex}
\begin{tikzpicture}
    \begin{axis}[...]
        \addplot [
            ...,
            error bars,
            x dir=both,
            x explicit,
            y dir=both,
            y explicit,
        ] table [
            x error expr=10,
            y error=PowerError,
        ] {resources/battery-title.csv};
    \end{axis}
\end{tikzpicture}
\end{minted}
\end{onlyenv}

\begin{onlyenv}<2->
\centering
\pgfplotstableset{
    col sep=comma,
}
\begin{tikzpicture}
    \begin{axis}[
        width=0.75\paperwidth,
        height=0.75\paperheight,
        xlabel={Resistencia [\unit{\ohm}]},
        ylabel={Potencia [\unit{\watt}]},
    ]
        \addplot [
            red,
            mark=*,
            only marks,
            error bars,
            x dir=both,
            x explicit,
            y dir=both,
            y explicit,
        ] table [
            x error expr=10,
            y error=PowerError,
        ] {resources/battery-title.csv};
    \end{axis}
\end{tikzpicture}
\end{onlyenv}
\end{frame}






\begin{frame}[fragile]
\begin{onlyenv}<1>
\begin{minted}{latex}
\begin{tikzpicture}
    \begin{axis}[...]
        \addplot [
            ...,
            error bars,
            x dir=both,
            x explicit,
            y dir=both,
            y explicit,
        ] table [
            x error index=2,
            y error=PowerError,
        ] {resources/battery-title.csv};
    \end{axis}
\end{tikzpicture}
\end{minted}
\end{onlyenv}

\begin{onlyenv}<2->
\centering
\pgfplotstableset{
    col sep=comma,
}
\begin{tikzpicture}
    \begin{axis}[
        width=0.75\paperwidth,
        height=0.75\paperheight,
        xlabel={Resistencia [\unit{\ohm}]},
        ylabel={Potencia [\unit{\watt}]},
    ]
        \addplot [
            red,
            mark=*,
            only marks,
            error bars,
            x dir=both,
            x explicit,
            y dir=both,
            y explicit,
        ] table [
            x error index=2,
            y error=PowerError,
        ] {resources/battery-title.csv};
    \end{axis}
\end{tikzpicture}
\end{onlyenv}
\end{frame}






\begin{frame}[fragile]
\begin{onlyenv}<1>
\begin{minted}{latex}
\begin{tikzpicture}
    \begin{axis}[...]
        \addplot [
            ...,
            error bars,
            x dir=both,
            x explicit,
            y dir=both,
            y explicit,
        ] table [
            x error minus expr=1,
            x error plus=ResistanceError,
            y error=PowerError,
        ] {resources/battery-title.csv};
    \end{axis}
\end{tikzpicture}
\end{minted}
\end{onlyenv}

\begin{onlyenv}<2->
\centering
\pgfplotstableset{
    col sep=comma,
}
\begin{tikzpicture}
    \begin{axis}[
        width=0.75\paperwidth,
        height=0.75\paperheight,
        xlabel={Resistencia [\unit{\ohm}]},
        ylabel={Potencia [\unit{\watt}]},
    ]
        \addplot [
            red,
            mark=*,
            only marks,
            error bars,
            x dir=both,
            x explicit,
            y dir=both,
            y explicit,
        ] table [
            x error minus expr=1,
            x error plus=ResistanceError,
            y error=PowerError,
        ] {resources/battery-title.csv};
    \end{axis}
\end{tikzpicture}
\end{onlyenv}
\end{frame}





\begin{frame}[fragile]
\frametitle{Regresiones Lineales}

Con {\lmr\PGFPlots} podemos realizar regresiones lineales de una manera relativamente sencilla. Con el paquete \texttt{pgfplotstable}

\bigskip

\begin{minted}[escapeinside=||]{latex}
create col/linear regression={|\textit{config}|}
\end{minted}

\bigskip\pause

\begin{columns}
    \begin{column}{0.5\textwidth}
        \begin{alertblock}{IMPORTANTE}
            Aunque esto es posible hacerlo, es recomendable no utilizar estos comandos sobre un conjunto de datos muy grande o podremos acabar con la máxima memoria reservada de {\lmr\LaTeX}
        \end{alertblock}
    \end{column}
\end{columns}
\end{frame}






\begin{frame}[fragile]
\begin{onlyenv}<1>
\begin{minted}{latex}
\begin{tikzpicture}
    \begin{axis}[...]
        \addplot [
            ...,
            forget plot,
        ] table {resources/battery-title.csv};
        \addplot [...] table [
            y={create col/linear regression={x=Resistance, y=Power}}
        ] {resources/battery-title.csv};
        \addlegendentry{%
            $\pgfmathprintnumber{\pgfplotstableregressiona} \cdot x
            \pgfmathprintnumber[print sign]{\pgfplotstableregressionb}$}

    \end{axis}
\end{tikzpicture}
\end{minted}
\end{onlyenv}

\begin{onlyenv}<2->
\centering
\pgfplotstableset{
    col sep=comma,
}
\begin{tikzpicture}
    \begin{axis}[
        width=0.75\paperwidth,
        height=0.75\paperheight,
        xlabel={Resistencia [\unit{\ohm}]},
        ylabel={Potencia [\unit{\watt}]},
    ]
        \addplot [
            red,
            mark=*,
            only marks,
            forget plot,
        ] table {resources/battery-title.csv};

        \addplot [
            red,
        ] table [
            y={create col/linear regression={x=Resistance, y=Power}}
        ] {resources/battery-title.csv};
        \addlegendentry{%
            $\pgfmathprintnumber{\pgfplotstableregressiona} \cdot x
            \pgfmathprintnumber[print sign]{\pgfplotstableregressionb}$}
    \end{axis}
\end{tikzpicture}
\end{onlyenv}
\end{frame}





\begin{frame}[fragile]
\frametitle{Inclusión de tablas}

Con \texttt{pgfplotstable} podemos incluir tablas externas, de igual manera que con las gráficas y darles formato directamente en {\lmr\LaTeX}.

\bigskip\pause

\begin{minted}[escapeinside=||]{latex}
\pgfplotstabletypeset[|\textit{opciones}|]{|\textit{archivo}|}
\end{minted}
\end{frame}





\begin{frame}[fragile]
\begin{onlyenv}<1>
\begin{minted}{latex}
\pgfplotstabletypeset{resources/battery.tsv}
\end{minted}
\end{onlyenv}

\begin{onlyenv}<2->
\centering
\tiny
\pgfplotstabletypeset{resources/battery.tsv}
\par
\end{onlyenv}
\end{frame}





\begin{frame}[fragile]
\begin{onlyenv}<1>
\begin{minted}{latex}
\pgfplotstableset{
    every head row/.style={before row=\toprule,after row=\midrule},
    every last row/.style={after row=\bottomrule}
}
\pgfplotstabletypeset{resources/battery.tsv}
\end{minted}
\end{onlyenv}

\begin{onlyenv}<2->
\centering
\tiny
\pgfplotstableset{
    every head row/.style={before row=\toprule,after row=\midrule},
    every last row/.style={after row=\bottomrule}
}
\pgfplotstabletypeset{resources/battery.tsv}
\par
\end{onlyenv}
\end{frame}





\begin{frame}[allowframebreaks]
\frametitle{Referencias}
\nocite{*}
\printbibliography
\end{frame}
\end{document}
